\chapter{Linear continua and structural elements}\label{ch:elem}

\begin{tcolorbox}[colback=gray!5,colframe=gray!50,title={In this chapter}]
Everything in the later chapters is built from element matrices, so this chapter starts at the root: how a weak form becomes a sparse stiffness matrix and a mass matrix, and which element families \fea{} offers for scalar fields, plane and solid continua, beams, plates and shells. It does not re-derive shape functions. It looks instead at the places where an element can be wrong while looking plausible: locking, integration rules and the choice of mass matrix. Each is followed by a test against an independent reference: patch tests on distorted meshes, pure bending, a slender cantilever of a nearly incompressible material, beam frequencies against closed forms, and a plate against the shell that must reproduce it.\par\smallskip
After it you should be able to say how a weak form becomes a stiffness and a mass matrix, recognise locking or a wrong integration rule from a test with an independent reference, and choose between a consistent and a lumped mass matrix, knowing the frequency error of each.
\end{tcolorbox}

\section{From the weak form to an element matrix}\label{sec:el-weak}

The smallest complete example is the scalar element \anchor{fea_engine.elements.scalar.Quad4Laplace}, which has one unknown per node. For the equation $-\nabla\!\cdot(\kappa\nabla u)=q$ on a domain $\Omega$ with $u=0$ on a part of its boundary, the weak form is: find $u$ such that for every admissible $v$
\begin{equation}\label{eq:el-weak}
\int_\Omega\kappa\,\nabla v\cdot\nabla u\,d\Omega=\int_\Omega v\,q\,d\Omega .
\end{equation}
Inserting the interpolation $u\approx\sum_a N_a u_a$ and the same functions for $v$ turns \eqref{eq:el-weak} into $\bm K\bm u=\bm f$ with
\begin{equation}\label{eq:el-K}
\bm K_e=\int_{\Omega_e}\bm B^{T}\kappa\,\bm B\,d\Omega,\qquad \bm B=\nabla\bm N,\qquad
\bm f_e=\int_{\Omega_e}\bm N^{T}q\,d\Omega .
\end{equation}
The integrals are computed by Gauss quadrature on the reference element, with $d\Omega=\det\bm J\,d\xi$ and $\bm B$ mapped by the Jacobian $\bm J$. The elasticity elements have the same structure with a vector field: $\bm K_e=\int\bm B^{T}\bm D\bm B\,d\Omega$, where $\bm B$ is the strain-displacement matrix and $\bm D$ the elasticity matrix \citep{zienkiewicz2013fem,bathe2014fep}. Everything else in this chapter is a choice of $\bm N$, of $\bm D$, or of the quadrature rule in \eqref{eq:el-K}.

\tikzfigure[0.98\linewidth]{ch03_weak2matrix}{From the weak form to the matrices. The choices that decide whether an element is right are the interpolation, the material matrix and the quadrature rule.}{fig:el-weak2matrix}

\begin{algbox}{Element stiffness and mass by Gauss quadrature}
\begin{pseudo}
\Require node coordinates $\bm x_e$, material matrix $\bm D$ (or $\kappa$), density $\rho$, Gauss rule $\{\bm\xi_g,w_g\}$
\State $\bm K_e\gets\bm 0$;\ $\bm M_e\gets\bm 0$
\For{each Gauss point $g$}
  \State $\bm N,\ \partial\bm N/\partial\bm\xi\gets$ shape functions at $\bm\xi_g$
  \State $\bm J\gets(\partial\bm N/\partial\bm\xi)\,\bm x_e$;\ $\det\bm J$;\ $\bm B\gets\bm J^{-1}\partial\bm N/\partial\bm\xi$ mapped into strain--displacement form
  \State $\bm K_e\gets\bm K_e+w_g\det\bm J\;\bm B^T\bm D\bm B$
  \State $\bm M_e\gets\bm M_e+w_g\det\bm J\;\rho\bm N^T\bm N$
\EndFor
\If{lumped mass requested}
  \State $\bm M_e\gets$ HRZ lumping of $\bm M_e$ \Comment{diagonal, same total mass}
\EndIf
\Ensure $\bm K_e$, $\bm M_e$
\end{pseudo}
\end{algbox}

The rule evaluates $n_g$ integrands of $O(d^2n_{str})$ each, and it sets the rank of $\bm K_e$. It goes wrong when it is too low (hourglass modes, rank deficiency) or too high for the physics (locking, \Cref{sec:el-locking}), and when the Jacobian is not positive. The code is \anchor{fea_engine.elements.base.Element} (\code{stiffness}, \code{mass}) and \code{assemble_lumped_mass}; the checks are \Cref{sec:el-verif} (patch tests, beams) and \Cref{fig:el-mass} (mass choices).

\paragraph{The mass matrix: consistent or lumped.}
The same interpolation gives the consistent mass matrix $\bm M_e=\int\rho\bm N^{T}\bm N\,d\Omega$, which is sparse with the pattern of $\bm K$. A lumped mass matrix is diagonal: \fea{} provides element-by-element HRZ lumping (\code{assemble_lumped_mass}), which keeps the total mass of the element. The diagonal form is what the explicit time integration of Chapter~\ref{ch:dyn} needs, since it makes the solve for the accelerations trivial. The choice has a measurable effect on the eigenvalues. For the lowest eigenvalue of the unit-square membrane, $\lambda_1=2\pi^2$ for a clamped edge, \Cref{fig:el-mass} shows that the consistent mass overestimates it and the lumped mass underestimates it, both with the error falling like $h^{\cecMassConsOrder}$ and $h^{\cecMassLumpOrder}$. On the finest mesh the consistent mass is high by \cecMassConsFine{} and the lumped mass is low by \cecMassLumpFine{}. The signs are the usual ones for these two choices, and an average of the two is a classical way to reduce the error; this chapter does not test that.

\begin{figure}[ht]
\centering
\includegraphics[width=0.55\linewidth]{figures/ch03_mass.pdf}
\caption{Lowest eigenvalue of the clamped unit-square membrane (\code{Quad4Laplace}), consistent against lumped mass: the relative error in absolute value. The consistent-mass eigenvalue is above the exact one and the lumped-mass eigenvalue is below it. Produced by \texttt{scripts/ch03\_assets.py}.}\label{fig:el-mass}
\end{figure}

\section{The element families}\label{sec:el-families}

\Cref{tab:el-elements} lists the elements treated in this chapter with their nodes, degrees of freedom per node and the \code{gauss_order} attribute of the class. Plane and solid elements share one idea: a displacement-based isoparametric element with polynomial interpolation. The plane elements use $\bm D$ for plane stress or plane strain; the 3-D elements use the $6\times6$ isotropic $\bm D$. Linear triangles and tetrahedra have constant strain; the 4-node quadrilateral and the 8-node hexahedron have a bilinear or trilinear field and therefore represent bending badly (\Cref{sec:el-locking}).

\begin{table}[ht]
\centering\small
\begin{tabular}{llrrr}
\toprule
family & element & nodes & dofs per node & \code{gauss\_order} \\
\midrule
scalar & \code{Quad4Laplace} & 4 & 1 & 2 \\
plane & \code{Tri3PlaneStress} & 3 & 2 & 1 \\
plane & \code{Tri6PlaneStress} & 6 & 2 & 2 \\
plane & \code{Quad4PlaneStress} & 4 & 2 & 2 \\
plane & \code{Quad8PlaneStress} & 8 & 2 & 3 \\
solid & \code{Tet4Solid3D} & 4 & 3 & 1 \\
solid & \code{Tet10Solid3D} & 10 & 3 & 2 \\
solid & \code{Hex8Solid3D} & 8 & 3 & 2 \\
solid & \code{Hex8SolidBbar} & 8 & 3 & 2 \\
solid & \code{Hex20Solid3D} & 20 & 3 & 3 \\
beam & \code{Beam2DEulerBernoulli} & 2 & 2 & 1 \\
beam & \code{Beam2DReissner} & 2 & 3 & 2 \\
beam & \code{Beam3DEulerBernoulli} & 2 & 6 & 2 \\
plate/shell & \code{Quad4MindlinPlate} & 4 & 3 & 2 \\
plate/shell & \code{Shell4MITC} & 4 & 6 & 2 \\
\bottomrule
\end{tabular}

\caption{Elements of this chapter. \code{gauss_order} is the attribute of the class as it stands in tag \booktag; its meaning is the number of points per direction for quadrilaterals and hexahedra, and a simplex rule order for triangles and tetrahedra.}\label{tab:el-elements}
\end{table}

\paragraph{Structural elements.}
The beam elements use a one-dimensional mesh. \anchor{fea_engine.elements.beams.Beam2DEulerBernoulli} has two dofs per node (transverse displacement and rotation) and cubic Hermite interpolation. \anchor{fea_engine.elements.beams.Beam2DReissner} adds shear deformation, so it has an axial displacement as well, and \anchor{fea_engine.elements.beams3d.Beam3DEulerBernoulli} has six dofs per node. The corotational versions of these beams belong to Chapter~\ref{ch:nonlin}. \anchor{fea_engine.elements.plates.Quad4MindlinPlate} is a 4-node plate with deflection and two rotations per node, and \anchor{fea_engine.elements.shells.Shell4MITC} is a 4-node shell with membrane, bending and transverse-shear action and six dofs per node, with the assumed transverse shear strain of \citet{dvorkin1984}.

\section{Locking and how the elements avoid it}\label{sec:el-locking}

Locking is an element behaving too stiffly for a deformation mode that its interpolation cannot reproduce without also producing a parasitic strain. It is the one error in this chapter that grows when the problem gets \emph{easier} (a thinner beam, a more incompressible material), which is why it needs a test of its own.

\tikzfigure[0.98\linewidth]{ch03_locking}{The mechanism of shear locking in a bilinear element under bending (schematic). The exact bending mode curves the edges and has no shear strain; the bilinear interpolation keeps the edges straight, which produces a parasitic shear strain and an artificially stiff element.}{fig:el-lockmech}

\paragraph{Bending of continuum elements.}
A 4-node quadrilateral in bending cannot represent the curvature of a beam without shear strain, and the parasitic shear stiffens the beam. \Cref{fig:el-bending} loads a plane-stress beam by a pure end moment and plots the ratio of the computed tip deflection to the exact one (the exact solution is quadratic in the displacements, so it lies in the space of the quadratic elements). The linear elements converge from below: at 4 elements along the beam the ratio is \cecBendQuadFour{} for \code{Quad4PlaneStress} and \cecBendTriFour{} for \code{Tri3PlaneStress}, and at 32 elements it is \cecBendQuadThirtytwo{} and \cecBendTriThirtytwo{}. The observed order of the quadrilateral error is \cecBendQuadOrder, close to the expected second order. The quadratic elements \code{Quad8PlaneStress} and \code{Tri6PlaneStress} reproduce the exact solution up to rounding (\cecBendQuadEightErr{} and \cecBendTriSixErr{} for the tested meshes), as the theory says they must.

\begin{figure}[ht]
\centering
\includegraphics[width=0.6\linewidth]{figures/ch03_bending.pdf}
\caption{Pure bending of a plane-stress beam, tip deflection divided by the exact value, against the number of elements along the beam. The quadratic elements are exact on these meshes. The helper is the one used by \texttt{fea\_engine/tests/test\_benchmarks\_convergence.py}.}\label{fig:el-bending}
\end{figure}

\paragraph{Volumetric locking and mean-dilatation B-bar.}
A nearly incompressible material ($\nu\to0.5$) forces $\operatorname{tr}\bm\varepsilon\to0$ at every Gauss point, which a low-order element satisfies only by not deforming. The mean-dilatation (B-bar) element of \citet{hughes1980bbar}, \anchor{fea_engine.elements.solids.Hex8SolidBbar}, replaces the volumetric part of $\bm B$ by its average over the element, which relaxes the constraint to one condition per element. \Cref{tab:el-locking} bends a slender cantilever ($L/h=20$, square section) with a tip load and compares the tip deflection with the Euler--Bernoulli value. At $\nu=0.3$ the plain 8-node hexahedron gives \cecLockHexPlainA{} of the beam value (shear locking, as in two dimensions), and when $\nu\to0.4999$ it collapses to \cecLockHexPlainB{}, because volumetric locking is added. The B-bar element gives \cecLockHexBbarA{} at $\nu=0.3$ and \cecLockHexBbarB{} at $\nu=0.4999$: it removes the \emph{additional} locking caused by incompressibility and leaves the shear locking untouched, which is the claim of the method and nothing more. The 20-node hexahedron gives \cecLockHexTwentyA{} at $\nu=0.3$ and \cecLockHexTwentyB{} at $\nu=0.4999$ on a mesh with the same node spacing, so it locks little in shear but still shows some volumetric locking.

\begin{table}[ht]
\centering\small
\begin{tabular}{lrrrr}
\toprule
element & mesh & dofs & $\nu=0.3$ & $\nu=0.4999$ \\
\midrule
\code{Hex8Solid3D} & $20\times 4\times 4$ & 1575 & 0.711 & 0.177 \\
\code{Hex8SolidBbar} & $20\times 4\times 4$ & 1575 & 0.729 & 0.745 \\
\code{Hex20Solid3D} & $10\times 2\times 2$ & 963 & 0.989 & 0.905 \\
\bottomrule
\end{tabular}

\caption{Tip deflection of a slender cantilever divided by the Euler--Bernoulli value, three hexahedra and two Poisson ratios. The 20-node mesh has the node spacing of the 8-node mesh. A 3-D beam does not reach exactly 1 even for a perfect element (root and shear effects), so only the differences between the rows are significant.}\label{tab:el-locking}
\end{table}

\paragraph{Shear locking in the plate and its cure.}
\anchor{fea_engine.elements.plates.Quad4MindlinPlate} uses full $2\times2$ Gauss integration for the bending energy and a reduced one-point rule for the transverse shear energy (selective reduced integration, the default \code{'sri'}). With full integration for both, the element locks as the plate gets thin: \Cref{tab:el-plate} gives the tip deflection of a cantilever strip ($h/L=0.01$, $W/L=0.2$) divided by the beam-theory value $PL^3/(3D_pW)$ with the plate rigidity $D_p=Eh^3/(12(1-\nu^2))$. This reference is approximate (it ignores the anticlastic effect and shear), so the converged ratio is \cecPlateSriFine{} and not 1; what the table shows is the difference between the columns. With selective reduced integration the ratio is \cecPlateSriCoarse{} on a mesh of 4 elements along the length. With full integration it is \cecPlateFullCoarse{} on the same mesh and still \cecPlateFullFine{} on 64 elements. The shell \code{Shell4MITC}, which uses a different cure for locking (assumed shear strains tied at the element edges), converges to the same value: on the finest mesh the ratio of the shell to the plate is \cecPlateRatioFS. A flat shell must reproduce the plate, so this agreement is a check of both elements; the membrane response of the shell in this test was exactly zero, as it must be for a flat homogeneous single layer.

\begin{table}[ht]
\centering\small
\begin{tabular}{rrrr}
\toprule
elements along $L$ & \code{Quad4MindlinPlate} (SRI) & same, full integration & \code{Shell4MITC} \\
\midrule
4 & 1.059 & 0.005 & 1.051 \\
8 & 1.081 & 0.021 & 1.069 \\
16 & 1.077 & 0.081 & 1.075 \\
32 & 1.078 & 0.262 & 1.077 \\
64 & 1.078 & 0.602 & 1.078 \\
\bottomrule
\end{tabular}

\caption{Tip deflection of a flat cantilever strip divided by the beam-theory value with plate rigidity (an approximate reference), against the number of elements along the length (a quarter as many across).}\label{tab:el-plate}
\end{table}

\section{Verification}\label{sec:el-verif}

\paragraph{Patch tests.}
An element passes the patch test if a mesh of distorted elements reproduces a linear displacement field exactly when the boundary nodes are given the field, and recovers the constant stress. \Cref{tab:el-patch} gives the largest relative error of the displacement and of the stress for the \cecPatchCount{} continuum elements on distorted meshes; the worst value is \cecPatchWorst{}, which is rounding. The patch test is necessary but not sufficient for convergence, and it says nothing about locking, which is why the tests above exist. The B-bar hexahedron is not in this table: it does not support stress recovery in the nodal form the test uses, so its patch test (displacements only) is the one of \texttt{tests/test\_hex8\_bbar.py}, which passes together with the other three tests of that file (4 passed in the run for this chapter).

\begin{table}[ht]
\centering\small
\begin{tabular}{lrr}
\toprule
element & displacement error & stress error \\
\midrule
\code{Quad4PlaneStress} & $5.5\times10^{-16}$ & $5.2\times10^{-15}$ \\
\code{Quad8PlaneStress} & $1.0\times10^{-15}$ & $7.6\times10^{-15}$ \\
\code{Tri3PlaneStress} & $3.7\times10^{-16}$ & $4.9\times10^{-15}$ \\
\code{Tri6PlaneStress} & $1.8\times10^{-15}$ & $1.7\times10^{-14}$ \\
\code{Hex8Solid3D} & $1.8\times10^{-16}$ & $8.9\times10^{-16}$ \\
\code{Hex20Solid3D} & $1.2\times10^{-15}$ & $2.7\times10^{-15}$ \\
\code{Tet4Solid3D} & $3.6\times10^{-16}$ & $8.9\times10^{-16}$ \\
\code{Tet10Solid3D} & $1.2\times10^{-15}$ & $1.9\times10^{-15}$ \\
\bottomrule
\end{tabular}

\caption{Patch tests on distorted meshes. The test helper is \texttt{\_patch} of \texttt{test\_benchmarks\_convergence.py}.}\label{tab:el-patch}
\end{table}

\paragraph{Beams.}
The Euler--Bernoulli element is exact for an end load: one element reproduces the tip deflection $PL^3/(3EI)$ to \cecEBTipErr, because the cubic Hermite interpolation contains the exact solution. For vibration the cubic Hermite element converges with order \cecEBFreqOrder{} in the lowest cantilever frequency (\cecEBFreqHz{} Hz for the test beam), the expected fourth order. The Reissner beam approaches the Timoshenko tip deflection $PL^3/(3EI)+PL/(GA)$ with order \cecTimoOrder{} (\Cref{fig:el-beams}); shear contributes \cecTimoShearShare{}\% of the tip deflection for this slender test beam, so it is a check of the formulation more than a demonstration of shear effects.

\begin{figure}[ht]
\centering
\includegraphics[width=0.95\linewidth]{figures/ch03_beams.pdf}
\caption{Convergence of the first cantilever frequency of the Euler--Bernoulli beam (left) and of the tip deflection of the Reissner beam (right) with the element size, with the observed order.}\label{fig:el-beams}
\end{figure}

\begin{table}[ht]
\centering\small
\begin{tabularx}{\linewidth}{@{}>{\raggedright\arraybackslash}p{0.44\linewidth} X >{\centering\arraybackslash}p{0.11\linewidth}@{}}
\toprule
Concept & Code & Script calls \\
\midrule
scalar element, weak form, mass & \anchor{fea_engine.elements.scalar.Quad4Laplace} & yes \\
plane and solid elements & \anchor{fea_engine.elements.solids.Quad4PlaneStress}, \anchor{fea_engine.elements.solids.Quad8PlaneStress}, \anchor{fea_engine.elements.solids.Tri3PlaneStress}, \anchor{fea_engine.elements.solids.Tri6PlaneStress}, \anchor{fea_engine.elements.solids.Hex8Solid3D}, \anchor{fea_engine.elements.solids.Hex20Solid3D}, \anchor{fea_engine.elements.solids.Tet4Solid3D}, \anchor{fea_engine.elements.solids.Tet10Solid3D} & yes \\
B-bar hexahedron & \anchor{fea_engine.elements.solids.Hex8SolidBbar} & yes \\
beams & \anchor{fea_engine.elements.beams.Beam2DEulerBernoulli}, \anchor{fea_engine.elements.beams.Beam2DReissner}, \anchor{fea_engine.elements.beams3d.Beam3DEulerBernoulli} & yes \\
plate, shell & \anchor{fea_engine.elements.plates.Quad4MindlinPlate}, \anchor{fea_engine.elements.shells.Shell4MITC} & yes \\
convergence studies & \anchor{fea_engine.convergence.run_study}, \anchor{fea_engine.convergence.observed_order} & yes \\
\bottomrule
\end{tabularx}
\caption{Where the concepts of this chapter live in the repository (tag \booktag). The last column says whether the chapter's script \texttt{scripts/chNN\_assets.py} calls, by name, the code listed in the row: yes (every name), part, or no. Code reached only through another routine counts as no, so the column is a lower bound on what this book exercises.}\label{tab:el-map}
\end{table}

\begin{limitsbox}
\begin{itemize}
\item \textbf{Locking is not removed by the patch test.} All elements of \Cref{tab:el-patch} pass it, and the plain 8-node hexahedron still gives \cecLockHexPlainA{} of the beam value. For slender or nearly incompressible problems, choose the element by the tests of \Cref{sec:el-locking}, not by the patch test.
\item \textbf{B-bar addresses volumetric locking only.} It does not cure shear locking (\Cref{tab:el-locking}), and the 20-node element still showed some volumetric locking at $\nu=0.4999$ on the mesh used.
\item \textbf{Stress recovery} is not available for every element (for example \code{Hex8SolidBbar}); use \code{at='centroid'} where the error message says so.
\item \textbf{The plate reference is approximate.} \Cref{tab:el-plate} compares with a strip of plate rigidity, so only differences between its columns are claimed.
\item \textbf{Not covered:} orthotropic and layered materials, thermal and piezoelectric coupling, shells in large rotation (Chapter~\ref{ch:nonlin} treats the corotational and the director shell only in outline), and meshes of high aspect ratio or strong distortion beyond the patch-test meshes.
\end{itemize}
\end{limitsbox}

\section{Summary}\label{sec:el-summary}

Every element is $\int\bm B^{T}\bm D\bm B$ evaluated by Gauss quadrature, plus a mass matrix that is consistent or lumped; the consistent mass overestimates the membrane eigenvalue and the lumped mass underestimates it. All continuum elements pass the patch test on distorted meshes to rounding, but the low-order ones lock in bending, and nearly incompressible material locks the plain hexahedron to \cecLockHexPlainB{} of the beam value. B-bar fixes the incompressible part only, selective reduced integration keeps the Mindlin plate usable when thin, and the MITC shell agrees with the plate in the flat limit. The beams are exact (Euler--Bernoulli, end load) or converge at the expected orders.

\section*{Exercises}\label{sec:el-exercises}
\begin{enumerate}
\item \emph{(Implementation)} Repeat \Cref{fig:el-mass} with the average of the consistent and the lumped matrix, $\bm M=\tfrac12(\bm M_c+\bm M_l)$. Does the error of $\lambda_1$ cancel to leading order, and what order does it have?
\item \emph{(Implementation)} Run the locking test of \Cref{tab:el-locking} on three meshes, $10\times2\times2$, $20\times4\times4$ and $40\times8\times8$, for the plain and the B-bar hexahedron. Does shear locking go away with refinement, and at what rate?
\item \emph{(Implementation)} Use \code{integration='full'} for the plate on a thicker strip ($h/L=0.1$). Is the locking still visible? Compare with the Timoshenko value.
\item \emph{(Interpretation)} Replace the end load of the cantilever by a distributed load. Is one Euler--Bernoulli element still exact? Why not?
\item \emph{(Implementation)} Distort the patch-test mesh more strongly (move the interior nodes by half an element). Does any element in \Cref{tab:el-patch} fail?
\item \emph{(Derivation)} For a fixed--free bar of one two-node element, find the lowest eigenvalue with the consistent mass matrix $\tfrac{\rho AL}{6}\left[\begin{smallmatrix}2&1\\1&2\end{smallmatrix}\right]$ and with the lumped matrix, and show that their errors against the exact $(\pi/2)^2EA/(\rho AL^2)$ are about $+21.6\,\%$ and $-18.9\,\%$.
\end{enumerate}

\chapterrefs
